% ============================================================
%  CHAPTER 3 : CONSENSUS MECHANISMS AND SCALABILITY
% ============================================================
\chapter{Consensus Mechanisms and Scalability Engineering}
\label{ch:consensus}

\noindent\textbf{Declaration of Contributions:} This chapter was primarily authored by \textbf{Anmol Goklani}, who researched and wrote the sections on the Byzantine Generals Problem, types of consensus mechanisms (PoW, PoS, PBFT), and scalability engineering approaches (sharding, Layer~2, rollups, block size). All team members reviewed the technical content and verified the comparative tables.

\bigskip

The previous chapter explained \emph{how} Bitcoin records and orders transactions. This chapter addresses two deeper questions: \emph{how do independent nodes agree on what the truth is?}\ (consensus) and \emph{how can blockchain handle more transactions?}\ (scalability).

\section{The Consensus Problem}
\label{sec:consensus-problem}

In a decentralised blockchain network, there is no central server that decides which transactions are valid and in what order they occurred. Instead, hundreds or thousands of independent nodes---each with its own copy of the ledger---must collectively agree on a single, consistent state. This is the \textbf{consensus problem}.

The difficulty is compounded by the fact that some nodes may be faulty (crashed or unresponsive) or even \emph{malicious} (deliberately sending incorrect or contradictory information). The theoretical underpinning of this challenge is the \textbf{Byzantine Generals Problem}, formulated by Lamport, Shostak, and Pease in 1982~\cite{lamport1982byzantine}.

\begin{definition}[Consensus Mechanism]
\label{def:consensus}
A consensus mechanism is a protocol that enables a decentralised network of nodes to agree on a single, consistent state of the blockchain, ensuring that all honest participants validate the same transactions and maintain a common transaction history---without relying on any central authority.
\end{definition}

Formally, a consensus mechanism defines a state transition function:
\begin{equation}
    C(S_t, \text{Tx}) \to S_{t+1}
    \label{eq:state-transition}
\end{equation}
where $S_t$ is the current state, $\text{Tx}$ is a set of new transactions, and $S_{t+1}$ is the next agreed-upon state.

Any practical consensus mechanism must satisfy three key properties:
\begin{enumerate}
    \item \textbf{Consistency} (Agreement): All honest nodes see the same state.
    \item \textbf{Validity} (Integrity): Only correctly formed and properly signed transactions are accepted.
    \item \textbf{Fault Tolerance} (Liveness): The system continues to operate even in the presence of faulty or malicious nodes (Byzantine resilience).
\end{enumerate}

\section{Types of Consensus Mechanisms}
\label{sec:consensus-types}

Different blockchain systems use different consensus mechanisms, each making different trade-offs among security, efficiency, and decentralisation. We examine the three most important ones.

\subsection{Proof of Work (PoW)}
\label{sec:pow-consensus}

\Gls{pow} was described in detail in \Cref{sec:pow-detail} as the mechanism used by Bitcoin. To summarise:
\begin{itemize}
    \item Nodes (miners) compete to solve a cryptographic puzzle: find a nonce such that $H(\text{block} \| \text{nonce}) < \text{target}$.
    \item The first miner to find a valid solution broadcasts the block to the network.
    \item Other nodes verify the solution (a single hash computation) and accept the block if valid.
    \item The miner is rewarded with newly created cryptocurrency and transaction fees.
\end{itemize}

\textbf{Strengths:} PoW is battle-tested (Bitcoin has been running since 2009 without a successful 51\% attack) and provides strong security guarantees.

\textbf{Weaknesses:} It is extremely \emph{energy-intensive}. The Bitcoin network alone consumes more electricity than many countries. It also has low throughput---Bitcoin processes approximately 7 \gls{tps}, compared to Visa's approximately 50,000 \gls{tps}.

\subsection{Proof of Stake (PoS)}
\label{sec:pos}

\Gls{pos}~\cite{king2012ppcoin} replaces computational puzzles with economic incentives. Instead of miners competing with hardware, \emph{validators} are selected to propose and validate blocks based on the amount of cryptocurrency they have ``staked'' (locked up as collateral).

The probability of being selected as a validator is proportional to one's stake:
\begin{equation}
    P(\text{validator } i) \propto \text{stake}_i
    \label{eq:pos}
\end{equation}

If a validator acts dishonestly (e.g., tries to approve fraudulent transactions), their staked funds can be ``slashed'' (confiscated)---providing a direct financial deterrent against misbehaviour.

\textbf{Strengths:} PoS is vastly more energy-efficient than PoW (Ethereum reduced its energy consumption by approximately 99.95\% after switching from PoW to PoS in September 2022). It also enables higher throughput.

\textbf{Weaknesses:} PoS can lead to \emph{wealth centralisation}---the rich get richer, since those with more stake are more likely to be selected as validators and earn rewards.

\subsection{Practical Byzantine Fault Tolerance (PBFT)}
\label{sec:pbft}

\Gls{pbft}~\cite{castro1999pbft} is a consensus protocol designed for permissioned blockchain networks (where the set of nodes is known in advance). It works through multiple rounds of voting among nodes.

PBFT can tolerate up to $f$ faulty (including malicious) nodes out of $n$ total nodes, provided:
\begin{equation}
    n \geq 3f + 1
    \label{eq:pbft-threshold}
\end{equation}
Agreement is reached when at least $2f + 1$ nodes vote for the same state.

\textbf{Strengths:} PBFT provides \emph{deterministic finality}---once a block is confirmed, it is final (no possibility of reorganisation). It is also fast for small networks.

\textbf{Weaknesses:} PBFT does not scale well. The communication overhead grows as $O(n^2)$ with the number of nodes, making it impractical for large, open networks.

\subsection{Comparison of Consensus Mechanisms}

\Cref{tab:consensus-comparison} provides a side-by-side comparison of the three consensus mechanisms.

\begin{table}[htbp]
\centering
\renewcommand{\arraystretch}{1.3}
\begin{tabularx}{\textwidth}{l X X X}
\toprule
\textbf{Property} & \textbf{PoW} & \textbf{PoS} & \textbf{PBFT} \\
\midrule
Selection method    & Computational puzzle      & Stake-weighted random      & Voting rounds \\
Energy usage        & Very high                 & Very low                   & Low \\
Throughput          & Low ($\sim$7 TPS)         & Medium--High               & High (small networks) \\
Finality            & Probabilistic             & Probabilistic/Deterministic & Deterministic \\
Fault tolerance     & $> 50\%$ honest hashpower & $> 50\%$ honest stake      & $n \geq 3f+1$ \\
Scalability         & Global, open              & Global, open               & Limited (permissioned) \\
Example             & Bitcoin                   & Ethereum 2.0               & Hyperledger Fabric \\
\bottomrule
\end{tabularx}
\caption{Comparison of Proof of Work, Proof of Stake, and PBFT consensus mechanisms.}
\label{tab:consensus-comparison}
\end{table}

The choice of consensus mechanism reflects a fundamental trade-off:
\begin{equation}
    \text{Security} \longleftrightarrow \text{Scalability} \longleftrightarrow \text{Decentralisation}
    \label{eq:trilemma-preview}
\end{equation}
No existing mechanism optimises all three simultaneously. This trade-off is explored in depth in \Cref{ch:challenges}.

\section{Scalability Engineering}
\label{sec:scalability}

\subsection{The Scalability Problem}

In most blockchain systems, \emph{every node} must verify and store \emph{every transaction}. While this redundancy is what makes blockchains secure and trustworthy, it severely limits throughput. \Cref{tab:throughput} illustrates the gap.

\begin{table}[htbp]
\centering
\renewcommand{\arraystretch}{1.2}
\begin{tabular}{l r}
\toprule
\textbf{System} & \textbf{Transactions Per Second (TPS)} \\
\midrule
Bitcoin          & $\sim$ 7 \\
Ethereum (pre-merge) & $\sim$ 15--30 \\
Visa             & $\sim$ 50{,}000 \\
\bottomrule
\end{tabular}
\caption{Throughput comparison: blockchain networks vs.\ traditional payment systems~\cite{kohad2020scalability}.}
\label{tab:throughput}
\end{table}

The goal of scalability engineering is to increase throughput without compromising security or decentralisation. Four main approaches have been proposed.

\subsection{Sharding}
\label{sec:sharding}

Sharding divides the blockchain network into smaller groups called \textbf{shards}, each responsible for processing a subset of transactions in parallel~\cite{luu2016sharding}.

\begin{figure}[htbp]
\centering
\begin{tikzpicture}[
    shard/.style={rectangle, draw, thick, rounded corners=5pt, minimum width=2.5cm, minimum height=2.5cm, font=\small},
    node_s/.style={circle, draw, fill=blue!15, minimum size=0.35cm},
    >=Stealth
]

\node[shard, fill=red!5] (s1) at (0, 0) {};
\node[above=0.05cm of s1.north, font=\footnotesize\bfseries] {Shard 1};
\foreach \i in {1,...,4} {
    \node[node_s] at ($(s1.center) + ({90*\i}:0.7)$) {};
}

\node[shard, fill=green!5] (s2) at (4, 0) {};
\node[above=0.05cm of s2.north, font=\footnotesize\bfseries] {Shard 2};
\foreach \i in {1,...,4} {
    \node[node_s] at ($(s2.center) + ({90*\i}:0.7)$) {};
}

\node[shard, fill=blue!5] (s3) at (8, 0) {};
\node[above=0.05cm of s3.north, font=\footnotesize\bfseries] {Shard 3};
\foreach \i in {1,...,4} {
    \node[node_s] at ($(s3.center) + ({90*\i}:0.7)$) {};
}

\draw[<->, thick, dashed, gray] (s1) -- (s2) node[midway, above, font=\scriptsize]{cross-shard};
\draw[<->, thick, dashed, gray] (s2) -- (s3) node[midway, above, font=\scriptsize]{cross-shard};

\end{tikzpicture}
\caption{Sharding: the network is partitioned into independent shards that process transactions in parallel. Cross-shard communication enables transactions that span multiple shards.}
\label{fig:sharding}
\end{figure}

\textbf{Advantages:} Throughput scales roughly linearly with the number of shards.

\textbf{Challenges:} Each shard has fewer nodes, which can reduce security. Cross-shard transactions require additional coordination protocols, adding complexity.

\subsection{Layer-2 Solutions}
\label{sec:layer2}

Layer-2 solutions move transaction processing \emph{off the main chain} (Layer~1), settling only the final result on-chain~\cite{poon2016lightning}. The main blockchain serves as a ``court of last resort'' that can resolve disputes.

Examples include:
\begin{itemize}
    \item \textbf{Payment Channels} (e.g., Bitcoin's Lightning Network): Two parties open a channel, conduct many transactions between themselves off-chain, and only broadcast the net result to the blockchain.
    \item \textbf{State Channels}: A generalisation of payment channels to arbitrary state transitions.
    \item \textbf{Sidechains}: Independent blockchains connected to the main chain via a two-way bridge.
\end{itemize}

\textbf{Advantages:} Near-instant transactions with very low fees.

\textbf{Challenges:} Added architectural complexity; users must stay online (or use watchtowers) to prevent fraud.

\subsection{Rollups}
\label{sec:rollups}

Rollups batch hundreds of transactions into a single transaction that is submitted to the main chain~\cite{buterin2021rollups}. There are two varieties:
\begin{itemize}
    \item \textbf{Optimistic Rollups:} Assume transactions are valid by default; any observer can submit a fraud proof within a challenge period if they detect an invalid transaction.
    \item \textbf{Zero-Knowledge (ZK) Rollups:} Each batch includes a cryptographic proof (a \gls{zkp}) that the batch is valid. No challenge period is needed because validity is mathematically guaranteed.
\end{itemize}

\textbf{Advantages:} Significant throughput improvement while inheriting the security of the main chain.

\textbf{Challenges:} Optimistic rollups have a delay for the challenge period; ZK rollups require complex proof generation.

\subsection{Larger Block Size}

The simplest approach to scaling is to increase the maximum size of each block, allowing more transactions per block.

\textbf{Advantages:} Straightforward to implement; immediate throughput improvement.

\textbf{Challenges:} Larger blocks take longer to propagate through the network, increasing the risk of forks. They also require more storage and bandwidth, potentially excluding nodes with limited resources and reducing decentralisation.

\subsection{Summary of Scaling Approaches}

\begin{table}[htbp]
\centering
\renewcommand{\arraystretch}{1.3}
\begin{tabularx}{\textwidth}{l X X}
\toprule
\textbf{Approach} & \textbf{Key Advantage} & \textbf{Main Challenge} \\
\midrule
Sharding            & Parallel processing, linear scaling     & Reduced per-shard security \\
Layer-2 Solutions   & Off-chain speed and low fees             & Architectural complexity \\
Rollups             & Batched efficiency, inherits L1 security & Proof generation overhead \\
Larger Block Size   & Simple implementation                    & Centralisation risk \\
\bottomrule
\end{tabularx}
\caption{Summary of blockchain scalability approaches and their trade-offs.}
\label{tab:scaling-summary}
\end{table}
